\section{Discussion and Limitations}

The paper should be explicit about scope. The current dataset supports Tangut short-text translation, especially title-like strings, not general open-domain Tangut sentence translation. The new subset analysis strengthens that boundary: the main gains from \texttt{baseline3\_2\_multitask} are concentrated on reference titles with canonical suffix patterns, while the judge advantage is much smaller on the non-suffix subset. That scope is still meaningful: many historically important Tangut materials appear precisely in compact catalog, title, or scholastic forms, and low-resource methods for such data are valuable.

The paper should also be honest about absolute difficulty. Even the strongest learned system reaches only 2.22/5 under the fixed-key reference-aware judge, so the task is far from solved. The result is therefore not that Tangut title translation is easy, but that carefully structured supervision and carefully constrained preference optimization each help different parts of the problem.

Two limitations matter most for future work. First, the real dataset is extremely small, so the paper cannot claim stable behavior beyond the current narrow domain. Second, the new DPO result is still metric-sensitive: the sigmoid variant is best under the reference-aware judge, while the robust+weighted variant is best on raw chrF++ and benefits most from title-oriented cleanup. The most useful next experiments are therefore not ``more DPO'' in the abstract, but stricter gap filtering, title-aware stopping/decoding control, and training objectives that reward both adequacy and clean title form directly.
